% !TeX spellcheck = ru_RU
\section{Введение}
\thispagestyle{withCompileDate}

Фандомные вики-ресурсы содержат тысячи страниц, описывающих персонажей, локации, события и их взаимосвязи в вымышленных вселенных. Ручное извлечение связей между сущностями требует огромных временных затрат и не масштабируется. Существуют общие инструменты для парсинга вики и библиотеки для распознавания именованных сущностей (NER), однако они не предоставляют готовой архитектуры, позволяющей легко переключаться между разными источниками данных (фандомами) и сохранять результаты в структурированном виде для дальнейшего анализа.

Проблема заключается в отсутствии лёгкого в реализации и одновременно гибкого программного решения, которое бы:
\begin{itemize}
    \item собирало текстовые данные с вики-сайта выбранного фандома;
    \item автоматически выделяло в тексте персонажей (PERSON) и локации (LOC);
    \item присваивало каждой сущности уникальный идентификатор, единый для всех документов;
    \item сохраняло связи между документами и сущностями в реляционную базу данных;
    \item позволяло отвечать на типовые вопросы аналитики (топ персонажей локации, список локаций персонажа, часто совместно упоминаемые сущности, поиск документов).
\end{itemize}

В рамках данной работы была разработана система, удовлетворяющая перечисленным требованиям. Архитектура построена таким образом, что замена парсера одного фандома на другой не влияет на остальные компоненты (извлечение сущностей, работа с базой данных, аналитические запросы). Для извлечения сущностей использована предобученная NER-модель \texttt{spaCy} (русская модель \texttt{ru\_core\_news\_sm}). Хранение данных организовано в PostgreSQL, импорт выполнен через DataGrip с решением характерных проблем форматирования CSV. Для проверки работоспособности реализованы SQL-запросы, отвечающие на заданные контрольные вопросы.

В отчёте описываются выбранные технологии, архитектура решения, детали реализации парсера и модуля извлечения сущностей, процесс загрузки данных в базу, а также приводятся примеры аналитических запросов и их результаты.
